\chapter*{For Grownups}
\markboth{For Grownups}{For Grownups}
\phantomsection
\addcontentsline{toc}{chapter}{For Grownups}

This book is written for curious young programmers, but it is also meant to be
comfortable for parents, teachers, clubs, and anyone helping from the side of
the keyboard.

The programs are intentionally small. A short listing gives a child a real win:
they can type the whole thing, run it, and point to the lines that make it work.
The point is not to hide the machine. The point is to make the machine friendly
enough that a beginner wants to keep going.

\section*{A Good Session}
\phantomsection
\addcontentsline{toc}{section}{A Good Session}

A good session is not measured by how many pages were finished. A good session
might be:

\begin{itemize}
\item typing ten lines without rushing,
\item finding one mistake,
\item changing a color and seeing the result,
\item making a silly sound effect,
\item explaining what one variable does.
\end{itemize}

\section*{Helping Without Taking Over}
\phantomsection
\addcontentsline{toc}{section}{Helping Without Taking Over}

When something breaks, ask small questions first. What line did NovaBASIC
mention? What changed since the last run? Does the line have the same number of
quote marks on both sides? This keeps the keyboard in the learner's hands.

\begin{novabox}[title=Why BASIC?]
BASIC is immediate. A beginner can type a line, run it, and see a result
without learning a build system first. NovaBASIC keeps that direct feeling but
adds modern Nova VM hardware for graphics, sprites, sound, and storage.
\end{novabox}

The style of this book is old-fashioned on purpose: listings, line numbers,
experiments, and plain explanations. The goal is not nostalgia by itself. The
goal is a path from curiosity to control.
